\chapter{Transcendental non-Archimedean metrics}
\index{non-Archimedean metric!transcendental}\index{inverse system!singularity types}
\label{chap:transcendental-metrics}

A pseudoeffective class need not contain a positive-mass current.
Consequently, a non-Archimedean metric in that class cannot in general
be represented by one positive-mass test curve.  The remedy of
\cite{DXZ23} is to remember compatible test curves after every K\"ahler
perturbation.  The compatible components keep zero-volume boundary
classes in the theory.

There are three layers of data, and later constructions use all three:
the cohomological class $\{\theta\}$, the compatible component at every
K\"ahler perturbation, and the Legendre parameter inside each component.
The unperturbed class and the inverse-limit metric may have zero volume.
Every perturbed component remains a positive-mass $\mathcal I$-model test curve.  Thus the inverse limit does
not manufacture a positive-mass potential in the unperturbed class when
none exists.

Throughout this chapter $X$ is a connected compact K\"ahler manifold of
dimension $n$, and $\theta$ is a smooth closed real $(1,1)$-form
whose class is pseudoeffective.  Let $\operatorname{Kah}(X)$ denote
the set of K\"ahler forms on $X$, ordered by
\[
                  \omega\preceq\omega'
                  \quad\Longleftrightarrow\quad
                  \omega\geq\omega'
\tag{53.1}\label{eq:tna-kahler-order}
\]
as forms.  This is a directed set: a sufficiently small positive
multiple of a fixed K\"ahler form is an upper bound in the order
\eqref{eq:tna-kahler-order}.
More precisely, fix a K\"ahler form $\omega_0$.
Compactness of $X$ implies that for every K\"ahler form $\omega$
there is an integer $m\geq1$ with $m^{-1}\omega_0\leq\omega$.
Thus $(m^{-1}\omega_0)_{m\geq1}$ is an increasing cofinal sequence
for this reversed order.

\section{The projective system}
\label{sec:tna-projective-system}

For each $\omega\in\operatorname{Kah}(X)$, the class
$\{\theta+\omega\}$ is big.  \Cref{chap:rays-test-curves}
therefore supplies the positive-mass space
\[
             \PSH^{\NA}(X,\theta+\omega)_{>0}
\tag{53.2}\label{eq:tna-positive-component-space}
\]
of $\mathcal I$-model test curves.

If $\omega\geq\omega'$, define
\[
 \rho_{\omega,\omega'}:
 \PSH^{\NA}(X,\theta+\omega')_{>0}
 \longrightarrow
 \PSH^{\NA}(X,\theta+\omega)_{>0},
 \qquad
 \rho_{\omega,\omega'}(\Gamma)
 :=P_{\theta+\omega}[\Gamma]_{\mathcal I},
\tag{53.3}\label{eq:tna-transition}
\]
where the envelope is taken level by level.

The difference $\omega-\omega'$ may be only semipositive.
The proof of \Cref{prop:rays-kahler-transition} still applies:
adding a smooth closed semipositive form preserves psh potentials,
and nonpluripolar multilinearity preserves positive mass.
Monotonicity and concavity of the $\mathcal I$-envelope then give
the same test-curve construction.  A zero difference gives the identity
on $\mathcal I$-model curves.

\begin{lemma}\label{lem:tna-transition-cocycle}
If $\omega\geq\omega'\geq\omega''$, then
\[
 \rho_{\omega,\omega'}\circ\rho_{\omega',\omega''}
 =\rho_{\omega,\omega''}.
\tag{53.4}\label{eq:tna-transition-cocycle}
\]
\end{lemma}

\begin{proof}
At every $\tau$ below the common terminal parameter, both sides are
$\mathcal I$-model potentials in $\PSH(X,\theta+\omega)$ and are
$\mathcal I$-equivalent to the same original level.  The uniqueness of
the normalized $\mathcal I$-model representative proves equality.
More explicitly, the first projection preserves every integer multiplier
ideal, and applying the second projection therefore imposes the same
multiplier-ideal constraints as projecting directly to
$\theta+\omega$.  Idempotence of the $\mathcal I$-envelope gives
\eqref{eq:tna-transition-cocycle} level by level.
The $-\infty$ levels agree by increasing continuity of the
$\mathcal I$-envelope.
\end{proof}

\begin{lemma}\label{lem:tna-transition-order}
Every transition $\rho_{\omega,\omega'}$ with $\omega\geq\omega'$
preserves and reflects order. In particular, it is injective.
\end{lemma}

\begin{proof}
Write $\alpha=\theta+\omega'$ and $\beta=\theta+\omega$.
For two $\mathcal I$-model potentials $u,v\in\PSH(X,\alpha)$,
the multiplier-ideal characterization of the envelope gives
\[
 \begin{split}
 u\leq v
 &\quad\Longleftrightarrow\quad
 u\preceq_{\mathcal I}v\\
 &\quad\Longleftrightarrow\quad
 P_\beta[u]_{\mathcal I}\leq P_\beta[v]_{\mathcal I}.
 \end{split}
\tag{53.4a}\label{eq:tna-transition-order}
\]
Here $u\preceq_{\mathcal I}v$ means
$\mathcal I(mu)\subseteq\mathcal I(mv)$ for every positive integer
$m$, independently of the background form.
The forward implication for envelopes follows from the candidate
description with $\preceq_{\mathcal I}$ in place of
$\sim_{\mathcal I}$; see
\cite[Corollary~6.1.3 and Proposition~6.1.3]{XiaBook}.
The reverse implication uses
$u\sim_{\mathcal I}P_\beta[u]_{\mathcal I}$ and
$v\sim_{\mathcal I}P_\beta[v]_{\mathcal I}$ from
\Cref{thm:complex-i-projection}, followed by the same candidate
description in the source background $\alpha$.
Thus monotonicity alone is not the order-reflection argument.

Apply \eqref{eq:tna-transition-order} at every common interior
parameter. Projection preserves the terminal parameter, so it
also detects whether the first curve has a longer effective interval
than the second. The closed endpoint convention then recovers
the comparison at every real parameter.
Reflecting order in both directions proves injectivity.
\end{proof}


\begin{definition}[Transcendental non-Archimedean metric]
\label{def:tna-pshna}
Define
\[
 \PSH^{\NA}(X,\theta)
 :=
 \varprojlim_{\omega\in\operatorname{Kah}(X)}
       \PSH^{\NA}(X,\theta+\omega)_{>0}
\tag{53.5}\label{eq:tna-inverse-limit}
\]
in the category of sets.  Thus
$\Gamma\in\PSH^{\NA}(X,\theta)$ is a family
\[
        \Gamma=(\Gamma^{\theta+\omega})_\omega,
\qquad
        \Gamma^{\theta+\omega}
          \in\PSH^{\NA}(X,\theta+\omega)_{>0},
\tag{53.6}\label{eq:tna-components}
\]
satisfying
\[
 P_{\theta+\omega}
 [\Gamma^{\theta+\omega'}]_{\mathcal I}
 =\Gamma^{\theta+\omega}
 \quad\text{whenever }\omega\geq\omega'.
\tag{53.7}\label{eq:tna-component-compatibility}
\]
\end{definition}

We order these families componentwise:
\[
 \Gamma\leq\Gamma'
 \quad\Longleftrightarrow\quad
 \Gamma^{\theta+\omega}_\tau\leq
 (\Gamma')^{\theta+\omega}_\tau
 \quad\text{for every }\omega\text{ and every }\tau.
\tag{53.7a}\label{eq:tna-componentwise-order}
\]
The closed endpoint convention is understood in this comparison.

Each component is a test curve, hence itself a family
$\tau\mapsto\Gamma_\tau^{\theta+\omega}$.  The two indices have
different roles:
\[
 \omega\quad\text{is the K\"ahler perturbation},\qquad
 \tau\quad\text{is the Legendre parameter}.
\tag{53.8}\label{eq:tna-two-indices}
\]
Confusing them reverses both limiting procedures.

\section{Components and maximal representatives}
\label{sec:tna-components}

\begin{proposition}\label{prop:tna-common-terminal}
For $\Gamma\in\PSH^{\NA}(X,\theta)$, the number
$\Gamma^{\theta+\omega}_{\max}$ is independent of $\omega$.
It is denoted $\Gamma_{\max}$.
\end{proposition}

\begin{proof}
Let $c_\omega:=\Gamma^{\theta+\omega}_{\max}$.  A transition changes a
finite level by an envelope but never reparameterizes it, and it sends an
identically $-\infty$ level to $-\infty$.  Hence
$c_\omega=c_{\omega'}$ whenever $\omega\geq\omega'$.  Given arbitrary
$\omega_1,\omega_2$, choose a K\"ahler form $\eta$ with
$\eta\leq\omega_1$ and $\eta\leq\omega_2$, for example a sufficiently
small multiple of a fixed K\"ahler form.  Compatibility compares both
components with the $\eta$-component and gives
$c_{\omega_1}=c_\eta=c_{\omega_2}$.
\end{proof}

\begin{corollary}[One component detects order]
\label{cor:tna-component-order}
For a fixed K\"ahler form $\omega_0$, evaluation
$\Gamma\mapsto\Gamma^{\theta+\omega_0}$ preserves and reflects order.
In particular, one component determines a compatible family uniquely.
\end{corollary}

\begin{proof}
Preservation is the definition \eqref{eq:tna-componentwise-order}.
Suppose the $\omega_0$-components are ordered, and fix any $\omega$.
Choose K\"ahler $\eta\leq\omega,\omega_0$.
By compatibility and \Cref{lem:tna-transition-order}, the comparison
at $\omega_0$ reflects to $\eta$ and then preserves to $\omega$.
This proves the comparison at every component.  Apply it in both
directions for uniqueness.
\end{proof}

This corollary is a uniqueness assertion.  Existence still requires a
compatible extension to every smaller perturbation; an arbitrary
positive-mass test curve at one perturbation need not provide one.

For fixed $\omega$ and $\tau<\Gamma_{\max}$, the level
$\Gamma_\tau^{\theta+\omega}$ is the largest normalized
$(\theta+\omega)$-psh potential in its $\mathcal I$-equivalence
class:
\[
 \Gamma_\tau^{\theta+\omega}
 =P_{\theta+\omega}
     [\Gamma_\tau^{\theta+\omega}]_{\mathcal I}.
\tag{53.9}\label{eq:tna-maximal-level}
\]
This is the maximal representative used in the construction.  It should
not be confused with a pointwise largest element of the whole space
$\PSH^{\NA}(X,\theta)$, which need not exist without a normalization.

The raw big-line-bundle filtration construction satisfies a corresponding
maximality statement by \Cref{thm:pseff-filtration-maximality}.  That theorem
is not a formal consequence of the levelwise normalization in
\eqref{eq:tna-maximal-level}: its proof resolves a fixed birational linear
system, quantizes the resulting obstacle, and returns complete sections to
the original graded series.

The $-\infty$ level is also compatible:
\[
 P_{\theta+\omega}
 [\Gamma_{-\infty}^{\theta+\omega'}]
 =
 P_{\theta+\omega}
 [\Gamma_{-\infty}^{\theta+\omega'}]_{\mathcal I}
 =
 \Gamma_{-\infty}^{\theta+\omega}
\tag{53.10}\label{eq:tna-minus-infinity-transition}
\]
for $\omega\geq\omega'$.  The first equality uses the
$\mathcal I$-goodness of the test-curve levels; it is not a general
identity for arbitrary quasi-psh functions.  Compatibility at this initial
level follows by increasing continuity of both envelopes; this is
\cite[Proposition~13.1.1]{XiaBook}.

\begin{proposition}\label{prop:tna-background-independence}
If $f\in C^\infty(X,\mathbf R)$ and $\theta'=\theta+\ddc f$, then
\[
                  \PSH^{\NA}(X,\theta)
                  \xrightarrow{\sim}
                  \PSH^{\NA}(X,\theta')
\tag{53.11}\label{eq:tna-background-independence}
\]
canonically, by replacing each level by the normalized envelope of its
translate by $-f$.  These identifications satisfy the cocycle
condition.
\end{proposition}

\begin{proof}
At each perturbed class this is \cite[Proposition~9.3.2]{XiaBook}.
Explicitly, replace a level $u$ by
$P_{\theta'+\omega}[u-f]_{\mathcal I}$.
The translated potential is $(\theta'+\omega)$-psh and has the same
multiplier ideals as $u$, since multiplication by $\exp(mf)$ is locally
bounded above and below.  Its mass before projection is unchanged
because the current is unchanged.  It remains $\mathcal I$-good under
this change of background, so its projected level has positive mass.
The resulting normalized $\mathcal I$-model representative is determined
by its ideal data in the new background.  Both orders of applying a
transition and this map therefore give the same representative.
Applying the map for $-f$ is its inverse, and successive changes by
$f$ and $g$ compose to the change by $f+g$.  Taking inverse limits
proves the assertion.  The envelopes themselves need not be pointwise
unchanged by a smooth translate.
\end{proof}

\section{Volume}
\label{sec:tna-volume}

The definition below is a limit of a net, so it needs the net to be
monotone.  The required identity is the one comparison that every mass
computation in this chapter performs, and it is isolated here rather than
repeated inside individual proofs.

\begin{lemma}[Transition mass identity]
\label{lem:tna-transition-mass-identity}
Let $\Gamma\in\PSH^{\NA}(X,\theta)$, let $\omega\geq\omega'$ be
K\"ahler forms, and let $\sigma$ be either $-\infty$ or a parameter
$\tau<\Gamma_{\max}$.  Put
\[
 S^\sigma_{\omega'}
 :=\theta+\omega'+\ddc\Gamma_\sigma^{\theta+\omega'},
 \qquad
 \gamma:=\omega-\omega'\geq0 .
\]
Then
\[
 \int_X\left(\theta+\omega+
       \ddc\Gamma_\sigma^{\theta+\omega}\right)^n
 =\int_X\left(S^\sigma_{\omega'}+\gamma\right)^n .
\tag{53.12}\label{eq:tna-transition-mass-identity}
\]
\end{lemma}

\begin{proof}
Compatibility \eqref{eq:tna-component-compatibility} at the finite levels,
and \eqref{eq:tna-minus-infinity-transition} at the $-\infty$ level, give
$\Gamma_\sigma^{\theta+\omega}
   =P_{\theta+\omega}[\Gamma_\sigma^{\theta+\omega'}]_{\mathcal I}$.
The left side of \eqref{eq:tna-transition-mass-identity} is therefore
$\vol(\theta+\omega,\Gamma_\sigma^{\theta+\omega'})$ in the sense of
\eqref{eq:complex-current-volume}.  The potential $u:=\Gamma_\sigma^{\theta+\omega'}$ is the normalized
$\mathcal I$-model representative in its own class, and positive-mass $\mathcal I$-model initial and interior levels are
$\mathcal I$-good, independently of the background, by
\cite[Example~7.1.2 and Theorem~7.1.1]{XiaBook}.  Hence
$P_{\theta+\omega}[u]=P_{\theta+\omega}[u]_{\mathcal I}$; this is the
identity already used at \eqref{eq:tna-minus-infinity-transition}.  The mass
identity \eqref{eq:complex-p-projection} for the $P$-envelope then gives
\[
 M_{\theta+\omega}(u)
 =M_{\theta+\omega}\bigl(P_{\theta+\omega}[u]\bigr)
 =M_{\theta+\omega}\bigl(P_{\theta+\omega}[u]_{\mathcal I}\bigr)
 =\vol(\theta+\omega,u),
\]
the equality case of \eqref{eq:complex-mass-volume}.  The
$\mathcal I$-goodness is essential and must not be replaced by an appeal to
\eqref{eq:complex-current-volume}: that formula supplies only
$M_{\theta+\omega}(u)\leq\vol(\theta+\omega,u)$, and it is emphatically
\emph{not} the case that the $\mathcal I$-envelope preserves mass --- see
the paragraph ``Not claimed'' in \Cref{sec:ch51-not-claimed} and
\Cref{ex:ch51-cantor}, where it raises the mass from $\tfrac12$ to $1$.
Expanding $\int_X(S^\sigma_{\omega'}+\gamma)^n$ gives the right side.
\end{proof}

The identity is used only through the two consequences recorded next.  It is
a statement about $\mathcal I$-model levels of one compatible family, not
about arbitrary quasi-psh potentials, and it is not asserted for a
perturbation $\gamma$ that fails to be semipositive.

\begin{definition}[Volume of a transcendental metric]
\label{def:tna-volume}
For $\Gamma\in\PSH^{\NA}(X,\theta)$, define
\[
 \vol(\Gamma)
 :=\lim_{\omega\in\operatorname{Kah}(X)}
 \int_X
 \left(\theta+\omega+
       \ddc\Gamma_{-\infty}^{\theta+\omega}\right)^n
 \in[0,+\infty).
\tag{53.13}\label{eq:tna-volume}
\]
\end{definition}

\begin{proposition}[The volume is well defined]
\label{prop:tna-volume-welldefined}
The net in \eqref{eq:tna-volume} is decreasing with respect to
\eqref{eq:tna-kahler-order} and bounded below by $0$.  Its limit
therefore exists and lies in $[0,+\infty)$, so
\Cref{def:tna-volume} is well posed.
\end{proposition}

\begin{proof}
Let $\omega\geq\omega'$ be K\"ahler, so that $\omega\preceq\omega'$ in
\eqref{eq:tna-kahler-order} and $\gamma=\omega-\omega'$ is smooth,
closed and semipositive.  Apply
\Cref{lem:tna-transition-mass-identity} with $\sigma=-\infty$ and expand
the right side of \eqref{eq:tna-transition-mass-identity} by multilinearity
of the nonpluripolar product:
\[
 \int_X\left(\theta+\omega+
       \ddc\Gamma_{-\infty}^{\theta+\omega}\right)^n
 =\sum_{a=0}^n\binom na
   \int_X\left(S^{-\infty}_{\omega'}\right)^a\wedge\gamma^{\,n-a}.
\]
Every summand is nonnegative, because $S^{-\infty}_{\omega'}$ is a
positive current and $\gamma$ is a semipositive smooth form.  Discarding
all summands except $a=n$ gives
\[
 \int_X\left(\theta+\omega+
       \ddc\Gamma_{-\infty}^{\theta+\omega}\right)^n
 \geq
 \int_X\left(\theta+\omega'+
       \ddc\Gamma_{-\infty}^{\theta+\omega'}\right)^n ,
\]
which is precisely decrease along \eqref{eq:tna-kahler-order}.  The net is
a net of nonpluripolar masses of positive currents, hence nonnegative.  A
decreasing net of real numbers bounded below converges to its infimum, and
the infimum lies in $[0,+\infty)$.
\end{proof}

The proof shows slightly more than monotonicity: the deficit between two
values of the net is the sum of the mixed terms with $a<n$, so two comparable perturbations have equal masses precisely when all
these mixed terms vanish.  This equality criterion does not require a
positive lower bound for the limiting volume.

\begin{remark}[Cofinal one-parameter computation]
\label{rem:tna-cofinal-volume}
Fix a K\"ahler form $\omega_0$.  The sequence
$(m^{-1}\omega_0)_{m\geq1}$ is cofinal, so
\[
 \vol(\Gamma)
 =\lim_{m\to\infty}
 \int_X
 \left(\theta+m^{-1}\omega_0+
       \ddc\Gamma_{-\infty}^{\theta+m^{-1}\omega_0}\right)^n.
\tag{53.14}\label{eq:tna-cofinal-volume}
\]
The definition remains a net definition; this formula is a convenient
cofinal computation, not a replacement of the index category in
functorial arguments.
Indeed, for every K\"ahler form $\omega$, compactness of $X$ gives
$m^{-1}\omega_0\leq\omega$ for all sufficiently large $m$.  This is
precisely cofinality in the reversed order
\eqref{eq:tna-kahler-order}.
\end{remark}

\section{The positive-volume regime}
\label{sec:tna-positive-volume}

There is already a positive-mass space
$\PSH^{\NA}(X,\theta)_{>0}$ when $\{\theta\}$ is big.  It embeds in
the inverse limit by
\[
 \Gamma\longmapsto
 \bigl(P_{\theta+\omega}[\Gamma]_{\mathcal I}\bigr)_\omega .
\tag{53.15}\label{eq:tna-positive-embedding}
\]

\begin{lemma}\label{lem:tna-mass-drop-transition}
Let $\Gamma\in\PSH^{\NA}(X,\theta)$, let $\omega$ be K\"ahler, and
let $0<t\leq1$.  For every $\tau<\Gamma_{\max}$,
\[
\begin{aligned}
&\int_X\left(\theta+t\omega+
       \ddc\Gamma_{-\infty}^{\theta+t\omega}\right)^n
 -\int_X\left(\theta+t\omega+
       \ddc\Gamma_{\tau}^{\theta+t\omega}\right)^n\\
&\quad\leq
\int_X\left(\theta+\omega+
       \ddc\Gamma_{-\infty}^{\theta+\omega}\right)^n
 -\int_X\left(\theta+\omega+
       \ddc\Gamma_{\tau}^{\theta+\omega}\right)^n .
\end{aligned}
\]
\end{lemma}

\begin{proof}
For $\sigma\in\{-\infty,\tau\}$, put
\[
 S_t^\sigma
 :=\theta+t\omega+\ddc\Gamma_\sigma^{\theta+t\omega}.
\]
This is the notation of
\Cref{lem:tna-transition-mass-identity} with $\omega'=t\omega$ and
$\gamma=(1-t)\omega$, which is semipositive because $t\leq1$.  That
lemma gives
\[
 \int_X\left(\theta+\omega+
       \ddc\Gamma_\sigma^{\theta+\omega}\right)^n
 =\int_X\left(S_t^\sigma+(1-t)\omega\right)^n.
\]
Expand the right side by multilinearity of the non-pluripolar product and
subtract the identities for $\sigma=-\infty$ and $\sigma=\tau$.  The
result is
\[
\sum_{a=0}^n\binom na(1-t)^{n-a}
\left(
 \int_X(S_t^{-\infty})^a\wedge\omega^{n-a}
 -\int_X(S_t^\tau)^a\wedge\omega^{n-a}
\right).
\]
Since
$\Gamma_{-\infty}^{\theta+t\omega}
 \geq\Gamma_\tau^{\theta+t\omega}$, every parenthesis is nonnegative by
mixed-mass monotonicity.  Keeping only the summand $a=n$ proves the
inequality.
\end{proof}

\begin{theorem}\label{thm:tna-positive-volume}
The image of \eqref{eq:tna-positive-embedding} is exactly
\[
       \{\Gamma\in\PSH^{\NA}(X,\theta)\mid\vol(\Gamma)>0\}.
\tag{53.16}\label{eq:tna-positive-volume-locus}
\]
For an object in this image,
\[
                 \vol(\Gamma)
                 =\int_X\theta_{\Gamma_{-\infty}}^n.
\tag{53.17}\label{eq:tna-positive-volume-formula}
\]
\end{theorem}

\begin{proof}
First let $\Lambda$ be a positive-mass $\mathcal I$-model test curve
in the unperturbed class, and put $u=\Lambda_{-\infty}$.
Its initial level is positive-mass and $\mathcal I$-good.
The argument of \Cref{lem:tna-transition-mass-identity} therefore gives
\[
 M_{\theta+m^{-1}\omega_0}
   \bigl(P_{\theta+m^{-1}\omega_0}[u]_{\mathcal I}\bigr)
 =\int_X(\theta_u+m^{-1}\omega_0)^n
 \longrightarrow M_\theta(u)>0.
\]
The convergence follows by expanding the finite mixed-mass polynomial.
Cofinality proves \eqref{eq:tna-positive-volume-formula} for the embedded
curve.  The embedding is injective by the same ideal-data argument as
\Cref{lem:tna-transition-order}.

Conversely, let $\vol(\Gamma)=\epsilon>0$ and fix $\omega_0$.
For every $\sigma<\Gamma_{\max}$ put
$v_{m,\sigma}=\Gamma_\sigma^{\theta+m^{-1}\omega_0}$.
Compatibility makes $v_{m,\sigma}$ decrease with $m$: an ideal envelope
is at least its normalized input.  Every level has supremum zero.
Their compact nested zero sets have a common point, so the decreasing
limit is not identically $-\infty$.  Each tail is psh in a fixed
background $\theta+m_0^{-1}\omega_0$; letting $m_0\to\infty$ shows
that the limit is $\theta$-psh.  Define
\[
 \Gamma'_\sigma:=\lim_{m\to\infty}v_{m,\sigma}.
\tag{53.18}\label{eq:tna-positive-recovery}
\]

We use precisely \cite[Proposition~3.1.12]{XiaBook}, equivalently
\cite[Lemma~2.10]{XiaOperationsV3}: for decreasing normalized
$P$-fixed potentials in backgrounds $\theta+\epsilon_m\omega_0$,
$\epsilon_m\downarrow0$, the limit is $P$-fixed and its mass is the
limit of their masses, including when that mass is zero.
When the limiting mass is positive, their divisorial Lelong numbers
also converge to those of the limit.  Here all approximants are
$P$-fixed by \Cref{lem:rays-levelwise-i-envelope}.  In particular,
\[
 M_\theta(\Gamma'_\sigma)
 =\lim_m M_{\theta+m^{-1}\omega_0}(v_{m,\sigma}).
\tag{53.18a}\label{eq:tna-recovered-mass}
\]
This analytic theorem, rather than a liminf estimate alone, identifies
the mass of the recovered potential.

Choose $\tau<\Gamma_{\max}$ sufficiently negative that
\[
 \begin{split}
 &M_{\theta+\omega_0}(\Gamma_{-\infty}^{\theta+\omega_0})
 -M_{\theta+\omega_0}(\Gamma_\tau^{\theta+\omega_0})
 <\epsilon/2.
 \end{split}
\tag{53.19}\label{eq:tna-mass-drop-choice}
\]
This follows from increasing continuity of model masses as
$\tau\to-\infty$, as in \Cref{lem:rays-test-positive-levels}.
By \Cref{lem:tna-mass-drop-transition} the drop at $m^{-1}\omega_0$
is no larger.  Since every initial component mass is at least its
infimum $\epsilon$, \eqref{eq:tna-recovered-mass} gives
$M_\theta(\Gamma'_\tau)\geq\epsilon/2$.
For $\sigma\leq\tau$ monotonicity gives the same lower bound.
If $\tau<\sigma<\Gamma_{\max}$, choose
$\sigma'\in(\sigma,\Gamma_{\max})$ and $a\in(0,1)$ with
$\sigma=a\tau+(1-a)\sigma'$.
Concavity, mixed-mass monotonicity and nonpluripolar multilinearity give
\[
 M_{\theta+m^{-1}\omega_0}(v_{m,\sigma})
 \geq a^nM_{\theta+m^{-1}\omega_0}(v_{m,\tau}).
\]
Consequently every interior $\Gamma'_\sigma$ has positive mass.

The limits preserve decrease and concavity in the parameter.
They are ideal-model as well: monotonicity of ideal envelopes and
inclusion of their background competitor sets give
\[
 \Gamma'_\sigma\leq P_\theta[\Gamma'_\sigma]_{\mathcal I}
 \leq P_{\theta+m^{-1}\omega_0}[\Gamma'_\sigma]_{\mathcal I}
 \leq P_{\theta+m^{-1}\omega_0}[v_{m,\sigma}]_{\mathcal I}
 =v_{m,\sigma}
\]
for every $m$, and taking the infimum proves equality on the left.
The usc initial level is positive-mass, being less singular than every
interior level, and is ideal-model by increasing ideal-envelope
continuity in \Cref{lem:rays-levelwise-i-envelope}.
Complete the curve at its terminal parameter by the closed endpoint
convention and set it equal to $-\infty$ above that parameter.
Thus $\Gamma'$ is an $\mathcal I$-model test curve.

It remains to recover the original components.  For a fixed interior
$\sigma$, compatibility says that all $v_{m,\sigma}$ are
$\mathcal I$-equivalent.  Their divisorial Lelong numbers are therefore
constant in $m$.  The positive-mass part of the analytic theorem above
and \eqref{eq:tna-recovered-mass} show that $\Gamma'_\sigma$ has these
same numbers.  Uniqueness of normalized ideal-model representatives
then yields
$P_{\theta+m^{-1}\omega_0}[\Gamma'_\sigma]_{\mathcal I}=v_{m,\sigma}$.
For an arbitrary K\"ahler $\omega$, choose
$m^{-1}\omega_0\leq\omega$ and apply the transition cocycle to obtain
$P_{\theta+\omega}[\Gamma'_\sigma]_{\mathcal I}
 =\Gamma_\sigma^{\theta+\omega}$.
Increasing continuity recovers the initial level, and the closed
convention recovers the endpoint.  Hence the embedding of $\Gamma'$
is exactly $\Gamma$.  Applying the first part of the proof to $\Gamma'$
also proves \eqref{eq:tna-positive-volume-formula} for this family.
The recovery theorem is \cite[Lemma~13.1.2]{XiaBook} and
\cite[Lemma~3.10]{XiaOperationsV3}.
\end{proof}

\begin{corollary}\label{cor:tna-positive-no-ambiguity}
When $\vol(\Gamma)>0$, one may work either with the recovered test
curve $\Gamma'$ or with the inverse system
$(\Gamma^{\theta+\omega})_\omega$.  The recovered curve is unique,
because equality after one K\"ahler transition identifies all
$\mathcal I$-model levels.
\end{corollary}

\begin{proof}
Existence of the recovered curve and recovery of the original compatible
family are the two conclusions of the preceding theorem.  If two
recovered curves give the same family, their levels have the same image
after a fixed K\"ahler transition.  The transition is injective on
positive-mass $\mathcal I$-model levels, so the curves agree at every
parameter.
\end{proof}

\section{The zero-volume regime}
\label{sec:tna-zero-volume}

If $\vol(\Gamma)=0$, construction
\eqref{eq:tna-positive-recovery} may produce levels of zero mass and need
not produce a test curve in the sense of
\Cref{def:rays-test-curve}.  The inverse system
\eqref{eq:tna-inverse-limit} is then the object.

\begin{remark}[Zero-volume firewall]
\label{rem:tna-zero-volume-firewall}
In zero volume we do not:
\begin{enumerate}
\item write a nonexistent unperturbed component
      $\Gamma^{\theta}$;
\item apply the positive-mass idempotence theorem for $P_\theta$;
\item infer a common positive lower bound for level masses;
\item subtract values of two levels at points where both are
      $-\infty$.
\end{enumerate}
Comparisons are made after a K\"ahler perturbation.  The volume is
computed by the decreasing net in \eqref{eq:tna-volume}.
\end{remark}

\begin{example}[A boundary class]
\label{ex:tna-boundary-class}
Let $\{\theta\}$ be pseudoeffective with $\vol\{\theta\}=0$.
Choose a normalized potential $V_\theta$ with minimal singularities in the
class $\{\theta\}$.  At every perturbation $\omega$, take the
constant-in-$\tau$ curve with level
$P_{\theta+\omega}[V_\theta]_{\mathcal I}$.  The transition cocycle
\eqref{eq:tna-transition-cocycle} makes these components a compatible
inverse system.  By \Cref{prop:tna-homogeneous-volume} its volume is
$\vol(\theta,V_\theta)=0$: every $\theta$-psh potential, including
$P_\theta[V_\theta]_{\mathcal I}$, has zero top nonpluripolar mass in
this nonbig class by \Cref{prop:ch51-mass-inequalities}.
This is the
homogeneous metric introduced explicitly in \Cref{def:tna-homogeneous}.
No positive-mass test curve in the unperturbed class can represent it,
since such a curve would make \eqref{eq:tna-positive-volume-formula}
positive.
\end{example}

\begin{example}[A concrete zero-volume class]
\label{ex:tna-product-zero-volume}
Let $X=C_1\times C_2$ be a product of compact Riemann surfaces and let
$\theta=\operatorname{pr}_1^*\omega_1$, where $\omega_1$ is K\"ahler
on $C_1$.  Then $\{\theta\}$ is nef and nonzero, but
$\int_X\theta^2=0$.  For each K\"ahler form $\omega$ on $X$, the
form $\theta+\omega$ is K\"ahler, and
$P_{\theta+\omega}[0]_{\mathcal I}=0$: the zero potential is an
admissible competitor and the obstacle bounds every competitor by zero.
Thus the component is the constant zero-level $\mathcal I$-model test
curve with endpoint zero.  The transition maps
give a compatible inverse system, and along a cofinal ray its volume is
\[
             \lim_{m\to\infty}\int_X
             (\theta+m^{-1}\omega)^2=\int_X\theta^2=0.
\]
Every component therefore has positive mass although the inverse-limit
object has volume zero.  This is the basic geometric reason the
perturbation index cannot be suppressed.
\end{example}

\begin{example}[Zero volume from a rigid negative curve]
\label{ex:tna-rigid-negative-curve}
The preceding class is nef.  Zero volume also occurs, for an unrelated
reason, at classes that are not nef, and the inverse system absorbs both
cases without modification.

Let $X$ be a compact K\"ahler surface carrying a smooth rational curve
$E$ with $E^2=-1$, for instance the blow-up of $\mathbf P^2$ at one
point, and let $\theta$ be a smooth closed real $(1,1)$-form with
$\{\theta\}=c_1(\mathcal O_X(E))$.  Choose a smooth Hermitian metric $h$ on $\mathcal O_X(E)$ with
curvature $\theta$, which is possible by the $\partial\bar\partial$-lemma.
Let $s_E$ be its canonical section, and normalize $\varphi=\log|s_E|_h^2$ so that
$\theta+\ddc\varphi=[E]$, the integration current --- the square is the
one demanded by the Lelong--Poincar\'e formula
$\ddc\log|s_D|_h^2=[D]-\Theta_D$ in the convention
$\ddc=\tfrac{i}{2\pi}\partial\bar\partial$; with $\log|s_E|_h$ one
would get $\tfrac12(\theta+[E])$ instead.

The class is pseudoeffective, being represented by an effective divisor, and
is not nef, since $\theta\cdot E=E^2=-1<0$.  The current $[E]$ is carried
by $E$, which is pluripolar, and the nonpluripolar product puts no mass on
a pluripolar set; hence $M_\theta(\varphi)=0$.  The potential $\varphi$
has analytic singularities, so \eqref{eq:complex-mass-volume} is an equality
and $\vol(\theta,\varphi)=0$.  By
\Cref{prop:tna-homogeneous-volume} the associated homogeneous metric has
\[
 \vol(\Gamma^\varphi)=\vol(\theta,\varphi)=0 .
\]

Its components are nevertheless of positive mass.  Fix a K\"ahler form
$\omega_0$ and put $\omega_m=m^{-1}\omega_0$.  Then
$\theta+\omega_m+\ddc\varphi=[E]+\omega_m$, which agrees with the smooth
form $\omega_m$ off $E$; since the nonpluripolar product ignores $E$,
\[
 \int_X\bigl(\theta+\omega_m+\ddc\varphi\bigr)^2
 =\int_X\omega_m^2
 =m^{-2}\int_X\omega_0^2>0 ,
\]
and this tends to $0$, which is the volume computed through the cofinal
formula \eqref{eq:tna-cofinal-volume}.

The two examples isolate the two mechanisms.  In
\Cref{ex:tna-product-zero-volume} the class is nef and its top
self-intersection vanishes.  Here the class is not nef, and its Zariski
decomposition is $\{E\}=0+\{E\}$: the positive part is zero and the class
is entirely negative part, which is why $\vol\{\theta\}=0$.  Neither
mechanism is detectable at a single perturbation, where the component always
has positive mass.
\end{example}

\section{Homogeneous metrics from complex singularities}
\label{sec:tna-homogeneous}

\begin{definition}[Homogeneous metric]
\label{def:tna-homogeneous}
Let $\varphi\in\PSH(X,\theta)$.  Define
$\Gamma^\varphi\in\PSH^{\NA}(X,\theta)$ by
\[
 \Gamma^\varphi_{\max}=0,\qquad
 \Gamma_{\tau}^{\varphi,\theta+\omega}
 =P_{\theta+\omega}[\varphi]_{\mathcal I}
 \quad(\tau<0).
\tag{53.20}\label{eq:tna-homogeneous}
\]
\end{definition}

This is well defined.  If $T=\theta+\ddc\varphi$, then
$T+\omega$ is positive and its top nonpluripolar mass is at least
$\int_X\omega^n>0$ by multilinearity: the smooth top term loses no
mass on the polar set.  Since
$P_{\theta+\omega}[\varphi]_{\mathcal I}\geq\varphi-\sup_X\varphi$,
its mass is positive too.  The projected level is $\mathcal I$-model
and hence fixed by $P$ by \Cref{lem:rays-levelwise-i-envelope}.
Constant interior levels, extended by their closed endpoint at zero,
therefore form a test curve.  The envelope cocycle makes these curves
compatible.  Thus positive mass of the original $T$ is unnecessary.


It is homogeneous in the sense that every negative level has the same
singularity representative.  Its Legendre transform at a perturbed
component is
\[
 (\Gamma^{\varphi,\theta+\omega})_t^*
       =P_{\theta+\omega}[\varphi]_{\mathcal I},
 \qquad t>0,
\tag{53.21}\label{eq:tna-homogeneous-legendre}
\]
with the prescribed value at $t=0$ given by the relevant
$-\infty$ model level.  The potential need not have positive
nonpluripolar mass in the unperturbed class.

\begin{proposition}\label{prop:tna-homogeneous-volume}
For every $\varphi\in\PSH(X,\theta)$,
\[
                \vol(\Gamma^\varphi)
                =\vol(\theta,\varphi),
\tag{53.22}\label{eq:tna-homogeneous-volume}
\]
where the right side is the multiplier-ideal volume
\eqref{eq:complex-current-volume}.  It need not equal the raw
nonpluripolar mass $M_\theta(\varphi)$.
\end{proposition}

\begin{proof}
Fix a K\"ahler form $\omega_0$ and set
$u_m=P_{\theta+m^{-1}\omega_0}[\varphi]_{\mathcal I}$.
The competitor families decrease as $m$ increases, so $u_m$ decreases.
Each $u_m$ is normalized and fixed by $P$; its limit $u$ is
$\theta$-psh and has supremum zero.  This follows either from compact
nested zero sets and decreasing psh limits, or directly from
\cite[Proposition~3.1.12]{XiaBook}.

We first identify $u$.  Each normalized competitor for
$P_\theta[\varphi]_{\mathcal I}$ is a competitor in every perturbed
background, so
$P_\theta[\varphi]_{\mathcal I}\leq u_m$ and hence
$P_\theta[\varphi]_{\mathcal I}\leq u$.
On the other hand $\varphi-\sup_X\varphi\leq u\leq u_m$.
For every prime divisor $E$ over $X$, monotonicity gives
\[
 \nu(\varphi,E)\geq\nu(u,E)\geq\nu(u_m,E)=\nu(\varphi,E).
\]
Thus $u\sim_{\mathcal I}\varphi$, by
\Cref{def:complex-i-preorder}, and $u\leq0$ makes it a competitor
for the unperturbed ideal envelope.  Consequently
$u=P_\theta[\varphi]_{\mathcal I}$.

The mass conclusion of \cite[Proposition~3.1.12]{XiaBook} now gives
\[
 \lim_{m\to\infty}M_{\theta+m^{-1}\omega_0}(u_m)
 =M_\theta(u)=\vol(\theta,\varphi).
\]
That conclusion requires the approximants to be fixed by $P$, which
has been verified, and holds even when the limiting mass is zero.
Cofinality identifies the left side with $\vol(\Gamma^\varphi)$.
The raw-mass distinction is \eqref{eq:complex-mass-volume}.
\end{proof}

\section{Bimeromorphic invariance}
\label{sec:tna-bimeromorphic}

Pullbacks of K\"ahler forms can be degenerate.  We therefore need the
following extension theorem before constructing the pullback map.

\begin{theorem}[Semipositive component extension; analytic input]
\label{thm:tna-semipositive-extension}
Let $\mathcal R(X)$ consist of the smooth closed real semipositive
$(1,1)$-forms $\eta$ with $\int_X\eta^n>0$, with the reversed form
order.  Restricting a compatible family to K\"ahler forms induces
\[
 \varprojlim_{\eta\in\mathcal R(X)}
       \PSH^{\NA}(X,\theta+\eta)_{>0}
 \xrightarrow{\sim}\PSH^{\NA}(X,\theta).
\tag{53.22a}\label{eq:tna-semipositive-extension}
\]
For any fixed K\"ahler form $\kappa$, its unique extension satisfies
\[
 \Gamma^{\theta+\eta}_\tau
 =\inf_{m\geq1}\Gamma^{\theta+\eta+m^{-1}\kappa}_\tau
 \quad(\tau<\Gamma_{\max}).
\tag{53.22b}\label{eq:tna-semipositive-component}
\]
The extended components have positive mass and the same ideal data
under the transition maps.
\end{theorem}

\begin{proof}[Source and scope of the input]
This is \cite[Proposition~13.1.5]{XiaBook}, with the explicit
construction and positive-mass conclusion in
\cite[Lemma~3.8, using Lemma~2.10]{XiaOperationsV3}.
It is imported as an analytic extension theorem.
The index category $\mathcal R(X)$ is not directed, as the footnote
to Proposition~13.1.5 specifies.  Thus
\eqref{eq:tna-semipositive-extension} is not an application of the
ordinary cofinal-net argument in \Cref{rem:tna-cofinal-volume}.
\end{proof}

Let $\pi:Y\to X$ be a proper bimeromorphic morphism from a compact
K\"ahler manifold.  For every K\"ahler form $\omega$ on $X$,
$\pi^*\omega\in\mathcal R(Y)$, since its top integral equals
$\int_X\omega^n>0$.

\begin{theorem}\label{thm:tna-bimeromorphic}
Pullback induces a canonical bijection
\[
 \pi^*:\PSH^{\NA}(X,\theta)
       \xrightarrow{\sim}
       \PSH^{\NA}(Y,\pi^*\theta).
\tag{53.23}\label{eq:tna-bimeromorphic}
\]
It preserves the common terminal parameter, volume, positivity of
volume, and homogeneous metrics:
\[
 \pi^*\Gamma^\varphi=\Gamma^{\pi^*\varphi}.
\tag{53.24}\label{eq:tna-bimeromorphic-homogeneous}
\]
\end{theorem}

\begin{proof}
First use the psh descent interface of
\cite[Section~4.8]{XiaOperationsV3} and the birational mass identity
\eqref{eq:complex-npp-modification}:
$\pi^*: \PSH(X,\alpha)\to\PSH(Y,\pi^*\alpha)$ is a bijection,
ideal equivalence is detected by the same prime divisors over either
space, and total nonpluripolar masses are preserved.
These are the levelwise inputs used in
\cite[Corollary~13.1.1]{XiaBook} and
\cite[Section~4.8]{XiaOperationsV3}.
Psh descent and surjectivity of $\pi$ preserve the zero obstacle,
so the defining competitor sets give
\[
 P_{\pi^*(\theta+\omega)}[\pi^*u]_{\mathcal I}
 =\pi^*P_{\theta+\omega}[u]_{\mathcal I}.
\tag{53.25}\label{eq:tna-level-bimeromorphic}
\]

For $\Gamma$ downstairs and a K\"ahler form $\eta$ upstairs,
choose a K\"ahler form $\omega$ downstairs so small that
$\pi^*\omega\leq\eta$.  Define the $\eta$-component levelwise by
\[
 (\pi^*\Gamma)^{\pi^*\theta+\eta}
 :=P_{\pi^*\theta+\eta}
       [\pi^*\Gamma^{\theta+\omega}]_{\mathcal I}.
\tag{53.25a}\label{eq:tna-pullback-component}
\]
Given two choices of $\omega$, choose a common smaller K\"ahler form
on $X$.  Equation~\eqref{eq:tna-level-bimeromorphic} and the envelope
cocycle show that both definitions agree.  The pulled-back level has
positive mass; adding the semipositive difference and projecting
preserves positivity.  The same cocycle proves compatibility for two
comparable $\eta$'s.  Hence this defines a metric upstairs.

Conversely, extend a metric $\Delta$ upstairs by
\Cref{thm:tna-semipositive-extension}.  For each $\omega$ downstairs,
descend its $\pi^*\omega$-component through the psh bijection.
Equation~\eqref{eq:tna-level-bimeromorphic} shows that the descended
curves are ideal-model and compatible, and mass invariance gives their
positive mass.  They define a metric downstairs.
The downstairs composite is the identity by levelwise pullback and
descent.  For the upstairs composite, choose $\pi^*\omega\leq\eta$;
compatibility of the extended family gives
\[
 P_{\pi^*\theta+\eta}
      [\Delta^{\pi^*\theta+\pi^*\omega}]_{\mathcal I}
 =\Delta^{\pi^*\theta+\eta}.
\]
This is exactly \eqref{eq:tna-pullback-component}, so the other composite
is the identity too.  The extension theorem supplies these semipositive
components; no cofinality of pullback forms among K\"ahler forms is
asserted.  All these operations preserve the terminal parameter.

To verify volume, fix a K\"ahler $\omega_0$ downstairs and let
$\Delta=\pi^*\Gamma$.
Write $A(\eta)$ for the initial mass of the extended $\eta$-component
upstairs.  For $\gamma=\pi^*(t\omega_0)$ and a K\"ahler $\kappa$
upstairs, compatibility and ideal-goodness of the positive initial
level give
\[
 A(\gamma+\delta\kappa)
 =\int_Y\bigl(\pi^*\theta+\gamma+
        \ddc\Delta_{-\infty}^{\pi^*\theta+\gamma}
        +\delta\kappa\bigr)^n
 \longrightarrow A(\gamma)
 \quad(\delta\downarrow0).
\]
Since $\gamma+\delta\kappa$ is K\"ahler, each of its masses is at
least $\vol(\Delta)$.  Thus $A(\gamma)\geq\vol(\Delta)$.
For any K\"ahler $\eta$ upstairs, sufficiently small $t$ satisfies
$\pi^*(t\omega_0)\leq\eta$, and the transition mass identity gives
$A(\pi^*(t\omega_0))\leq A(\eta)$.
It follows that
\[
 \vol(\Delta)
 =\lim_{t\downarrow0}A(\pi^*(t\omega_0)).
\]
The pulled-back components have the same initial masses as the
$t\omega_0$-components downstairs.  Cofinality downstairs therefore
identifies this limit with $\vol(\Gamma)$.
Finally \eqref{eq:tna-level-bimeromorphic} applied to the constant
negative levels proves \eqref{eq:tna-bimeromorphic-homogeneous}.
\end{proof}

\begin{remark}[Why smoothness is stated]
\label{rem:tna-smooth-scope}
The construction extends in \cite{Xia23Operations} to reduced,
irreducible, unibranch compact K\"ahler spaces.  This chapter keeps $X$
smooth because \Cref{chap:complex-prerequisites} supplied the
multiplier-ideal and nonpluripolar interfaces in that setting.
No normality, resolution-independence, or restriction statement on a
singular complex space is inferred merely from
\eqref{eq:tna-bimeromorphic}.
\end{remark}

\begin{remark}[Source note]
\label{rem:tna-source-note}
Definitions \eqref{eq:tna-inverse-limit} and
\eqref{eq:tna-volume}, the homogeneous embedding, and bimeromorphic
invariance are the transcendental construction of \cite{DXZ23}.  The
positive-volume recovery in
\Cref{thm:tna-positive-volume,cor:tna-positive-no-ambiguity} is
\cite[Lemma~3.10]{XiaOperationsV3}; the semipositive extension is
its Lemma~3.8.  The precise version checked here is the 22 October 2024
arXiv v3, corresponding to the subsequently published
\cite{Xia23Operations}.  The perturbation order,
$\mathcal I$-transition maps, and zero-volume safeguards have been
spelled out because they are the hypotheses used by every operation in
the next chapter.
\end{remark}
